% GENERATED FILE. Edit canonical lessons, then run npm run generate:books.
% canonical-source-hash: fnv1a64:f773c4319061d849
% canonical-lessons: GU-C465-sakhat, GU-C465-ciknu, GU-C465-uddhat, GU-C465-ahankar, GU-C465-tevavu

\chapter{Strict and Proud}
\label{ch:gu-strict-and-proud}
\hlchaptermodality{465}

\begin{chapteropening}
\textbf{By the end of this chapter:} \emph{I can describe a strict or proud person.}

Complete the last lesson of chapter 465: I can describe a strict or proud person.
\end{chapteropening}

\section[\texorpdfstring{\gu{સખત} (sakhat)}{સખત (sakhat)}]{\gu{સખત} (sakhat) — strict; severe}

\label{lesson:GU-C465-sakhat}



\begin{quote}
Before the new one: say the Gujarati for composed, then the Gujarati for to heed.
\end{quote}

\subsection*{You'll want to know: \gu{સખત}}
\textbf{\gu{સખત}} — \emph{sakhat} — \textquotedblleft{}strict; severe\textquotedblright{}.

\begin{quote}
\gu{અમારા} \gu{આચાર્ય} \gu{ઘણા} \gu{સખત} \gu{છે}.
\end{quote}

\begin{quote}
\emph{amārā ācārya ghaṇā sakhat che.}
\end{quote}

\begin{quote}
Our principal is very strict.
\end{quote}

\begin{cousinweb}
\textbf{\gu{સખત}} (\emph{sakhat}) is usually traced to Persian \emph{saxt}, \textquotedblleft{}hard, firm\textquotedblright{}.
\end{cousinweb}

\begin{grammarlens}[title={What you've built}]
One more word for describing what you lived and felt.
\end{grammarlens}

\subsection*{Your turn}
\begin{itemize}
\raggedright
  \item \emph{Say it:} \emph{sakhat}
  \item \emph{Say it:} \emph{sakhat}, once more
  \item \emph{Say it:} \emph{amārā ācārya ghaṇā sakhat che.} — Our principal is very strict.
  \item \emph{Recall:} say the Gujarati for composed, then the Gujarati for to heed, then say \emph{sakhat} again
\end{itemize}

\begin{tcolorbox}[breakable,colback=teal!4,colframe=teal!35!black,title={Before you move on}]
What does \gu{સખત} mean? (\textquotedblleft{}Strict; severe\textquotedblright{}.) Say it once more.
\end{tcolorbox}
\section[\texorpdfstring{\gu{ચીકણું} (cīkṇũ)}{ચીકણું (cīkṇũ)}]{\gu{ચીકણું} (cīkṇũ) — fussy; sticky}

\label{lesson:GU-C465-ciknu}



\begin{quote}
Before the new one: say the Gujarati for to heed, then the Gujarati for strict.
\end{quote}

\subsection*{You'll want to know: \gu{ચીકણું}}
\textbf{\gu{ચીકણું}} — \emph{cīkṇũ} — \textquotedblleft{}fussy; sticky\textquotedblright{}.

\begin{quote}
\gu{મારી} \gu{સાસુ} \gu{ઘરકામમાં} \gu{ઘણી} \gu{ચીકણી} \gu{છે}.
\end{quote}

\begin{quote}
\emph{mārī sāsu gharkāmmā̃ ghaṇī cīkṇī che.}
\end{quote}

\begin{quote}
My mother-in-law is very fussy about the housework.
\end{quote}

\begin{cousinweb}
\textbf{\gu{ચીકણું}} (\emph{cīkṇũ}) goes back to Sanskrit \emph{cikkaṇa}, \textquotedblleft{}slippery, unctuous\textquotedblright{}; CDIAL lists the Gujarati word, \textquotedblleft{}sticky\textquotedblright{} (CDIAL 4782), and Gujarati also uses it for a fussy person.
\end{cousinweb}

\begin{grammarlens}[title={What you've built}]
One more word for describing what you lived and felt.
\end{grammarlens}

\subsection*{Your turn}
\begin{itemize}
\raggedright
  \item \emph{Say it:} \emph{cīkṇũ}
  \item \emph{Say it:} \emph{cīkṇũ}, once more
  \item \emph{Say it:} \emph{mārī sāsu gharkāmmā̃ ghaṇī cīkṇī che.} — My mother-in-law is very fussy about the housework.
  \item \emph{Recall:} say the Gujarati for to heed, then the Gujarati for strict, then say \emph{cīkṇũ} again
\end{itemize}

\begin{tcolorbox}[breakable,colback=teal!4,colframe=teal!35!black,title={Before you move on}]
What does \gu{ચીકણું} mean? (\textquotedblleft{}Fussy; sticky\textquotedblright{}.) Say it once more.
\end{tcolorbox}
\section[\texorpdfstring{\gu{ઉદ્ધત} (uddhat)}{ઉદ્ધત (uddhat)}]{\gu{ઉદ્ધત} (uddhat) — insolent; rude}

\label{lesson:GU-C465-uddhat}



\begin{quote}
Before the new one: say the Gujarati for strict, then the Gujarati for fussy.
\end{quote}

\subsection*{You'll want to know: \gu{ઉદ્ધત}}
\textbf{\gu{ઉદ્ધત}} — \emph{uddhat} — \textquotedblleft{}insolent; rude\textquotedblright{}.

\begin{quote}
\gu{ઉદ્ધત} \gu{છોકરાએ} \gu{ગુરુને} \gu{ખરાબ} \gu{જવાબ} \gu{આપ્યો}.
\end{quote}

\begin{quote}
\emph{uddhat chokrāe gurune kharāb javāb āpyo.}
\end{quote}

\begin{quote}
The insolent boy gave the teacher a rude answer.
\end{quote}

\begin{cousinweb}
\textbf{\gu{ઉદ્ધત}} (\emph{uddhat}) is a tatsama, Sanskrit \emph{uddhata}, \textquotedblleft{}raised, lifted up\textquotedblright{}, as Monier-Williams records it; the sense \textquotedblleft{}haughty, insolent\textquotedblright{} grew from holding one's head too high.
\end{cousinweb}

\begin{grammarlens}[title={What you've built}]
One more word for describing what you lived and felt.
\end{grammarlens}

\subsection*{Your turn}
\begin{itemize}
\raggedright
  \item \emph{Say it:} \emph{uddhat}
  \item \emph{Say it:} \emph{uddhat}, once more
  \item \emph{Say it:} \emph{uddhat chokrāe gurune kharāb javāb āpyo.} — The insolent boy gave the teacher a rude answer.
  \item \emph{Recall:} say the Gujarati for strict, then the Gujarati for fussy, then say \emph{uddhat} again
\end{itemize}

\begin{tcolorbox}[breakable,colback=teal!4,colframe=teal!35!black,title={Before you move on}]
What does \gu{ઉદ્ધત} mean? (\textquotedblleft{}Insolent; rude\textquotedblright{}.) Say it once more.
\end{tcolorbox}
\section[\texorpdfstring{\gu{અહંકાર} (ahankār)}{અહંકાર (ahankār)}]{\gu{અહંકાર} (ahankār) — egotism; conceit}

\label{lesson:GU-C465-ahankar}



\begin{quote}
Before the new one: say the Gujarati for fussy, then the Gujarati for insolent.
\end{quote}

\subsection*{You'll want to know: \gu{અહંકાર}}
\textbf{\gu{અહંકાર}} — \emph{ahankār} — \textquotedblleft{}egotism; conceit\textquotedblright{}.

\begin{quote}
\gu{પૈસાનો} \gu{અહંકાર} \gu{સારો} \gu{નથી}.
\end{quote}

\begin{quote}
\emph{paisāno ahankār sāro nathī.}
\end{quote}

\begin{quote}
Pride in money is not good.
\end{quote}

A masculine noun.

\begin{cousinweb}
\textbf{\gu{અહંકાર}} (\emph{ahankār}) is a tatsama, Sanskrit \emph{ahaṃkāra}, the making of an \textquotedblleft{}I\textquotedblright{}, egotism, as Monier-Williams records it.
\end{cousinweb}

\begin{grammarlens}[title={What you've built}]
One more word for describing what you lived and felt.
\end{grammarlens}

\subsection*{Your turn}
\begin{itemize}
\raggedright
  \item \emph{Say it:} \emph{ahankār}
  \item \emph{Say it:} \emph{ahankār}, once more
  \item \emph{Say it:} \emph{paisāno ahankār sāro nathī.} — Pride in money is not good.
  \item \emph{Recall:} say the Gujarati for fussy, then the Gujarati for insolent, then say \emph{ahankār} again
\end{itemize}

\begin{tcolorbox}[breakable,colback=teal!4,colframe=teal!35!black,title={Before you move on}]
What does \gu{અહંકાર} mean? (\textquotedblleft{}Egotism; conceit\textquotedblright{}.) Say it once more.
\end{tcolorbox}
\section[\texorpdfstring{\gu{ટેવાવું} (ṭevāvũ)}{ટેવાવું (ṭevāvũ)}]{\gu{ટેવાવું} (ṭevāvũ) — to get used to}

\label{lesson:GU-C465-tevavu}



\begin{quote}
Before the new one: say the Gujarati for insolent, then the Gujarati for egotism.
\end{quote}

\subsection*{You'll want to know: \gu{ટેવાવું}}
\textbf{\gu{ટેવાવું}} — \emph{ṭevāvũ} — \textquotedblleft{}to get used to\textquotedblright{}.

\begin{quote}
\gu{નવા} \gu{ઘરમાં} \gu{બાળકો} \gu{જલદી} \gu{ટેવાઈ} \gu{ગયાં}.
\end{quote}

\begin{quote}
\emph{navā gharmā̃ bāḷko jaldī ṭevāī gayā̃.}
\end{quote}

\begin{quote}
The children quickly got used to the new house.
\end{quote}

Past (he …): \gu{તે} \gu{ટેવાઈ} \gu{ગયો}, \emph{te ṭevāī gayo}.

\begin{cousinweb}
\textbf{\gu{ટેવાવું}} (\emph{ṭevāvũ}) is built on \textbf{\gu{ટેવ}} (\emph{ṭev}), \textquotedblleft{}a habit\textquotedblright{}, which CDIAL traces to a reconstructed \emph{*ṭēv-}, \textquotedblleft{}appoint, fix\textquotedblright{} (CDIAL 5475), with the verb ending \emph{-āvũ}.
\end{cousinweb}

\begin{grammarlens}[title={What you've built}]
One more word for describing what you lived and felt.
\end{grammarlens}

\subsection*{Your turn}
\begin{itemize}
\raggedright
  \item \emph{Say it:} \emph{ṭevāvũ}
  \item \emph{Say it:} \emph{ṭevāvũ}, once more
  \item \emph{Say it:} \emph{navā gharmā̃ bāḷko jaldī ṭevāī gayā̃.} — The children quickly got used to the new house.
  \item \emph{Recall:} say the Gujarati for insolent, then the Gujarati for egotism, then say \emph{ṭevāvũ} again
\end{itemize}

\begin{tcolorbox}[breakable,colback=teal!4,colframe=teal!35!black,title={Before you move on}]
What does \gu{ટેવાવું} mean? (\textquotedblleft{}To get used to\textquotedblright{}.) Say it once more.
\end{tcolorbox}
